% GENERATED FILE. Edit canonical lessons, then run npm run generate:books.
% canonical-source-hash: fnv1a64:638f13886467637d
% canonical-lessons: JA-C521-shidaini, JA-C521-dandan, JA-C521-jojoni, JA-C521-youyaku, JA-C521-zuibun, JA-R521-review-516-521

\chapter{Little by Little}
\label{ch:ja-little-by-little}
\hlchaptermodality{521}

\begin{chapteropening}
\textbf{By the end of this chapter:} \emph{I can say how things will change over time.}

Complete the last lesson of chapter 521: I can say how things will change over time.
\end{chapteropening}

\section[\texorpdfstring{\ja{しだいに} (shidaini)}{しだいに (shidaini)}]{\ja{しだいに} (shidaini) — gradually}

\label{lesson:JA-C521-shidaini}



\begin{quote}
Before the new one: say the Japanese for beyond expectation, then the Japanese for sure.
\end{quote}

\subsection*{You'll want to know: \ja{しだいに}}
\textbf{\ja{しだいに}} — \emph{shidaini} — \textquotedblleft{}gradually\textquotedblright{}.

\begin{quote}
\textbf{\ja{はるに} \ja{なったら、しだいに} \ja{あたたかく} \ja{なるでしょう。}} — \emph{haru ni nattara, shidaini atatakaku naru deshō} — \textquotedblleft{}When spring comes, it will gradually get warmer.\textquotedblright{}
\end{quote}

\begin{cousinweb}
\textbf{\ja{しだいに}} (\emph{shidaini}) is Sino-Japanese (kango), with the particle \emph{ni}: two Chinese characters read with their on'yomi, readings borrowed from Chinese. They are \emph{shi}, \textquotedblleft{}next\textquotedblright{} (Middle Chinese \emph{tshijH}) and \emph{dai}, \textquotedblleft{}No.\textquotedblright{} (\emph{dejH}) (KANJIDIC2; Middle Chinese in Baxter's transcription).
\end{cousinweb}

\begin{grammarlens}[title={What you've built}]
One more word for saying what happens if.
\end{grammarlens}

\subsection*{Your turn}
\begin{itemize}
\raggedright
  \item \emph{Say it:} \emph{shidaini}
  \item \emph{Say it:} \emph{shidaini}, once more
  \item \emph{Say it:} \emph{haru ni nattara, shidaini atatakaku naru deshō} — When spring comes, it will gradually get warmer.
  \item \emph{Recall:} say the Japanese for beyond expectation, then the Japanese for sure, then say \emph{shidaini} again
\end{itemize}

\begin{tcolorbox}[breakable,colback=teal!4,colframe=teal!35!black,title={Before you move on}]
What does \ja{しだいに} mean? (\textquotedblleft{}Gradually\textquotedblright{}.) Say it once more.
\end{tcolorbox}
\section[\texorpdfstring{\ja{だんだん} (dandan)}{だんだん (dandan)}]{\ja{だんだん} (dandan) — little by little; gradually}

\label{lesson:JA-C521-dandan}



\begin{quote}
Before the new one: say the Japanese for sure, then the Japanese for gradually.
\end{quote}

\subsection*{You'll want to know: \ja{だんだん}}
\textbf{\ja{だんだん}} — \emph{dandan} — \textquotedblleft{}little by little; gradually\textquotedblright{}.

\begin{quote}
\textbf{\ja{くすりを} \ja{のんだら、ちちは} \ja{だんだん} \ja{よく} \ja{なるでしょう。}} — \emph{kusuri o nondara, chichi wa dandan yoku naru deshō} — \textquotedblleft{}If he takes the medicine, my father will gradually get better.\textquotedblright{}
\end{quote}

\begin{cousinweb}
\textbf{\ja{だんだん}} (\emph{dandan}) is Sino-Japanese (kango): one Chinese character read with its on'yomi, reading borrowed from Chinese. It is \emph{dan}, \textquotedblleft{}grade\textquotedblright{} (Middle Chinese \emph{dwanH}) (KANJIDIC2; Middle Chinese in Baxter's transcription).
\end{cousinweb}

\begin{grammarlens}[title={What you've built}]
One more word for saying what happens if.
\end{grammarlens}

\subsection*{Your turn}
\begin{itemize}
\raggedright
  \item \emph{Say it:} \emph{dandan}
  \item \emph{Say it:} \emph{dandan}, once more
  \item \emph{Say it:} \emph{kusuri o nondara, chichi wa dandan yoku naru deshō} — If he takes the medicine, my father will gradually get better.
  \item \emph{Recall:} say the Japanese for sure, then the Japanese for gradually, then say \emph{dandan} again
\end{itemize}

\begin{tcolorbox}[breakable,colback=teal!4,colframe=teal!35!black,title={Before you move on}]
What does \ja{だんだん} mean? (\textquotedblleft{}Little by little; gradually\textquotedblright{}.) Say it once more.
\end{tcolorbox}
\section[\texorpdfstring{\ja{じょじょに} (jojoni)}{じょじょに (jojoni)}]{\ja{じょじょに} (jojoni) — step by step; gradually}

\label{lesson:JA-C521-jojoni}



\begin{quote}
Before the new one: say the Japanese for gradually, then the Japanese for little by little.
\end{quote}

\subsection*{You'll want to know: \ja{じょじょに}}
\textbf{\ja{じょじょに}} — \emph{jojoni} — \textquotedblleft{}step by step; gradually\textquotedblright{}.

\begin{quote}
\textbf{\ja{おとうとは} \ja{れんしゅうしたら、じょじょに} \ja{はやく} \ja{なるでしょう。}} — \emph{otōto wa renshū shitara, jojoni hayaku naru deshō} — \textquotedblleft{}If my younger brother practises, he will slowly get faster.\textquotedblright{}
\end{quote}

\begin{cousinweb}
\textbf{\ja{じょじょに}} (\emph{jojoni}) is Sino-Japanese (kango), with the particle \emph{ni}: one Chinese character read with its on'yomi, reading borrowed from Chinese. It is \emph{jo}, \textquotedblleft{}gradually\textquotedblright{} (Middle Chinese \emph{zjo}) (KANJIDIC2; Middle Chinese in Baxter's transcription).
\end{cousinweb}

\begin{grammarlens}[title={What you've built}]
One more word for saying what happens if.
\end{grammarlens}

\subsection*{Your turn}
\begin{itemize}
\raggedright
  \item \emph{Say it:} \emph{jojoni}
  \item \emph{Say it:} \emph{jojoni}, once more
  \item \emph{Say it:} \emph{otōto wa renshū shitara, jojoni hayaku naru deshō} — If my younger brother practises, he will slowly get faster.
  \item \emph{Recall:} say the Japanese for gradually, then the Japanese for little by little, then say \emph{jojoni} again
\end{itemize}

\begin{tcolorbox}[breakable,colback=teal!4,colframe=teal!35!black,title={Before you move on}]
What does \ja{じょじょに} mean? (\textquotedblleft{}Step by step; gradually\textquotedblright{}.) Say it once more.
\end{tcolorbox}
\section[\texorpdfstring{\ja{ようやく} (yōyaku)}{ようやく (yōyaku)}]{\ja{ようやく} (yōyaku) — finally (after a long wait); at last}

\label{lesson:JA-C521-youyaku}



\begin{quote}
Before the new one: say the Japanese for little by little, then the Japanese for step by step.
\end{quote}

\subsection*{You'll want to know: \ja{ようやく}}
\textbf{\ja{ようやく}} — \emph{yōyaku} — \textquotedblleft{}finally (after a long wait); at last\textquotedblright{}.

\begin{quote}
\textbf{\ja{あと} \ja{いちにち} \ja{まったら、ようやく} \ja{やすみに} \ja{なるでしょう。}} — \emph{ato ichinichi mattara, yōyaku yasumi ni naru deshō} — \textquotedblleft{}If I wait one more day, I will finally have a holiday.\textquotedblright{}
\end{quote}

\begin{cousinweb}
\textbf{\ja{ようやく}} (\emph{yōyaku}) is a native Japanese word (wago), inherited rather than borrowed from Chinese, and usually written in kana; its older origin is not settled (JMdict).
\end{cousinweb}

\begin{grammarlens}[title={What you've built}]
One more word for saying what happens if.
\end{grammarlens}

\subsection*{Your turn}
\begin{itemize}
\raggedright
  \item \emph{Say it:} \emph{yōyaku}
  \item \emph{Say it:} \emph{yōyaku}, once more
  \item \emph{Say it:} \emph{ato ichinichi mattara, yōyaku yasumi ni naru deshō} — If I wait one more day, I will finally have a holiday.
  \item \emph{Recall:} say the Japanese for little by little, then the Japanese for step by step, then say \emph{yōyaku} again
\end{itemize}

\begin{tcolorbox}[breakable,colback=teal!4,colframe=teal!35!black,title={Before you move on}]
What does \ja{ようやく} mean? (\textquotedblleft{}Finally (after a long wait); at last\textquotedblright{}.) Say it once more.
\end{tcolorbox}
\section[\texorpdfstring{\ja{ずいぶん} (zuibun)}{ずいぶん (zuibun)}]{\ja{ずいぶん} (zuibun) — very; a good deal}

\label{lesson:JA-C521-zuibun}



\begin{quote}
Before the new one: say the Japanese for step by step, then the Japanese for finally (after a long wait).
\end{quote}

\subsection*{You'll want to know: \ja{ずいぶん}}
\textbf{\ja{ずいぶん}} — \emph{zuibun} — \textquotedblleft{}very; a good deal\textquotedblright{}.

\begin{quote}
\textbf{\ja{もし} \ja{まいにち} \ja{あるいたら、そふは} \ja{ずいぶん} \ja{つよく} \ja{なるでしょう。}} — \emph{moshi mainichi aruitara, sofu wa zuibun tsuyoku naru deshō} — \textquotedblleft{}If my grandfather walks every day, he will get a good deal stronger.\textquotedblright{}
\end{quote}

\begin{cousinweb}
\textbf{\ja{ずいぶん}} (\emph{zuibun}) is Sino-Japanese (kango): two Chinese characters read with their on'yomi, readings borrowed from Chinese, though JMdict marks it as usually written in kana. They are \emph{zui}, \textquotedblleft{}follow\textquotedblright{} (Middle Chinese \emph{zjwe}) and \emph{bun}, \textquotedblleft{}part\textquotedblright{} (\emph{bjunH}); how the two came to mean \textquotedblleft{}very\textquotedblright{} no dictionary we use records (KANJIDIC2, JMdict; Middle Chinese in Baxter's transcription).
\end{cousinweb}

\begin{grammarlens}[title={What you've built}]
One more word for saying what happens if.
\end{grammarlens}

\subsection*{Your turn}
\begin{itemize}
\raggedright
  \item \emph{Say it:} \emph{zuibun}
  \item \emph{Say it:} \emph{zuibun}, once more
  \item \emph{Say it:} \emph{moshi mainichi aruitara, sofu wa zuibun tsuyoku naru deshō} — If my grandfather walks every day, he will get a good deal stronger.
  \item \emph{Recall:} say the Japanese for step by step, then the Japanese for finally (after a long wait), then say \emph{zuibun} again
\end{itemize}

\begin{tcolorbox}[breakable,colback=teal!4,colframe=teal!35!black,title={Before you move on}]
What does \ja{ずいぶん} mean? (\textquotedblleft{}Very; a good deal\textquotedblright{}.) Say it once more.
\end{tcolorbox}
\section[(dialogue)]{Looking back: from Getting Ahead to Little by Little}

\label{lesson:JA-R521-review-516-521}



\begin{quote}
Say the Japanese for to get ahead (in life) and to get over.
\end{quote}

\subsection*{Your turn}
Say these aloud:

\begin{itemize}
\raggedright
  \item to get ahead (in life), to get over, all one's strength, the will to do something, to tackle (a challenge)
  \item to make a firm decision, to resolve (to do), to give up all thought of, to judge, an option
  \item to get back, to learn one's lesson, so close, to slip up, a fault
  \item to think again, to think better of (doing), to give up on (someone), to betray, safe (as a choice)
  \item a worry, helpless, unexpected, beyond expectation, sure
  \item gradually, little by little, step by step, finally (after a long wait), a good deal
\end{itemize}

\begin{tcolorbox}[breakable,colback=teal!4,colframe=teal!35!black,title={Before you move on}]
To get ahead (in life)? (\textbf{\ja{しゅっせする}}, \emph{shusse suru}.) To get over? (\textbf{\ja{のりこえる}}, \emph{norikoeru}.) All one's strength? (\textbf{\ja{ぜんりょく}}, \emph{zenryoku}.) The will to do something? (\textbf{\ja{やるき}}, \emph{yaruki}.) To tackle (a challenge)? (\textbf{\ja{いどむ}}, \emph{idomu}.) To make a firm decision? (\textbf{\ja{けつだんする}}, \emph{ketsudan suru}.) To resolve (to do)? (\textbf{\ja{けついする}}, \emph{ketsui suru}.) To give up all thought of? (\textbf{\ja{おもいきる}}, \emph{omoikiru}.) To judge? (\textbf{\ja{はんだんする}}, \emph{handan suru}.) An option? (\textbf{\ja{せんたくし}}, \emph{sentakushi}.) To get back? (\textbf{\ja{とりもどす}}, \emph{torimodosu}.) To learn one's lesson? (\textbf{\ja{こりる}}, \emph{koriru}.) So close? (\textbf{\ja{おしい}}, \emph{oshii}.) To slip up? (\textbf{\ja{しくじる}}, \emph{shikujiru}.) A fault? (\textbf{\ja{おちど}}, \emph{ochido}.) To think again? (\textbf{\ja{おもいなおす}}, \emph{omoinaosu}.) To think better of (doing)? (\textbf{\ja{おもいとどまる}}, \emph{omoitodomaru}.) To give up on (someone)? (\textbf{\ja{みかぎる}}, \emph{mikagiru}.) To betray? (\textbf{\ja{うらぎる}}, \emph{uragiru}.) Safe (as a choice)? (\textbf{\ja{ぶなん}}, \emph{bunan}.) A worry? (\textbf{\ja{きがかり}}, \emph{kigakari}.) Helpless? (\textbf{\ja{こころぼそい}}, \emph{kokorobosoi}.) Unexpected? (\textbf{\ja{おもいがけない}}, \emph{omoigakenai}.) Beyond expectation? (\textbf{\ja{よそうがい}}, \emph{yosōgai}.) Sure? (\textbf{\ja{かくじつ}}, \emph{kakujitsu}.) Gradually? (\textbf{\ja{しだいに}}, \emph{shidaini}.) Little by little? (\textbf{\ja{だんだん}}, \emph{dandan}.) Step by step? (\textbf{\ja{じょじょに}}, \emph{jojoni}.) Finally (after a long wait)? (\textbf{\ja{ようやく}}, \emph{yōyaku}.) A good deal? (\textbf{\ja{ずいぶん}}, \emph{zuibun}.)
\end{tcolorbox}
